\scanpage{232}
% corollary and question
Donc on trouve

\uline{Corollaire}
\[
 \mathrm{Diff}^{(i)}_+(E,P)\simeq
 \mathrm{Diff}^{i-1}(E\otimes\Omega^1_{S/T},P)/
 \operatorname{Im}\mathrm{Diff}^{i-2}(E\otimes\Omega^{\unsure{2}}_{S/T},P)
\]
% right edge of the preceding denominator meets the scan boundary
Par suite, on a un hom naturel injectif \textcircled{?}
\[
 \mathrm{Diff}^{(i)}_+(E,P)\hookrightarrow
 \mathrm{Hom}_{\mathrm{add}}(\operatorname{Ker}d_E^1,P)
\]
\inserted{l'image est appelée : faisceau des op. diff. de $\operatorname{Ker}d_E^1$ dans $P$}

% universal target
Pour $i$ fixé, le $P$ ``universel'' pour les op. diff à coeff constants (d'ordre $\leq i$) sur $E$ est $\unsure{p|P}^{(i)}(E)/\underline{\mathcal O}_S\cdot\delta_i^{\deleted{\unsure{E}}}(E)$, où
\[
 \delta_i:E\longrightarrow p^i(E)
\]
est défini par la stratification.

% 3), submodule
3) Supposons maintenant donné un sous-Module $N$ de $E$, on s'intéresse aux hom ``différentiels'' $\operatorname{Ker}d_E^1\longrightarrow P$ qui s'annulent sur $d_E^0 N$, qui correspondent donc aux \deleted{\unsure{hom}} op. diff. $E\longrightarrow P$, à coeff. cts, et qui s'annulent sur $N$, donc aux op. diff. $M=E/N\longrightarrow P$ qui ``\deleted{\unsure{fon}} \deleted{sont nuls sur}'' deviennent à coeff. cts sur $E$. Ils correspondent \inserted{(pour l'ordre $\leq i$)} aux hom linéaires $p^i(M)\longrightarrow P$ qui s'annulent sur $\rho_i E$ ; où
\[
 \rho_i:E\longrightarrow p^i(M)
\]
% archive sheet number: 231
